\documentclass{ltjsbook}
\usepackage{docmute} 
\usepackage{../myStyle} 

\begin{document} 
\chapter{導入；多変量データと心理学} 
\label{chapter1} 

この授業は基礎的な心理統計を修めた人向けの，応用コースになっています。基礎的な心理統計，という言葉に私が込めた意味は，
\begin{itemize}
    \item 記述統計；得られたデータの統計量を算出したり，可視化することによってデータの特徴を把握する。
    \item 推測統計；得られたデータが母集団からの標本であると考え，標本の特徴を使って母数を推測する。さらに標本の特徴から母集団について何らかの判断(意思決定)を行う。
\end{itemize}
という2点です。特に推測統計学の領域では，確率の話やさまざまな推定法，それに伴う技術などが必要ですから，これだけでも膨大な量だったのではと思います。
こうした基礎的な知識や技術を身につけると，とりあえず手元のデータを使って差があるとかないとか，どの程度効果があったのか，と言った基本的な判断はできると思います。複雑な計算式のところは機械(統計環境Rなど)がやってくれますので，出た結果だけをみて判断すれば良いのです\footnote{もちろん結果の見方が間違っていたり，拙速だったりしてはいけないなど，注意すべきところは多々あります。}。

しかし基礎的なところだけで満足していると，「はて，何がしたかったのかな」とそもそもの問題意識を忘れてしまうことにもなりかねません。とりあえず教わった(膨大な！)プロセスを経て実験結果を見てみればOKなんでしょう，というだけでは表面的な理解に止まっていると言わざるを得ません。さまざまなシーンに適用される方法論，その計算は何を意味していて，式にはどのような意味があり，式から得られるものは何をどこまで指し示しているのか，というところまで考えるのが，この講義の狙いです。数式は数式に過ぎない，というのはその通りなのですが，その数式にどのような意味があるのか，数式から得られる結果にどのような意味があるのか，を理解した上で数値(データ)を扱うようになることが目的です。数値だけに振り回される状態からの脱却，が狙いなわけです。

この授業の前半(前期)は，講義の形をとります。テーマとしては，\keyterm{因子分析}{いんしぶんせき}{Factor Analysis}を扱います。因子分析は，いわゆる質問紙調査を行った後で適用される多変量解析法の一種で，たくさんの質問項目の背後にある「因子」を見つけ出します。その因子にはXXX特性，XXX傾向といった名前がつけられ，その上で心理学的に考察することが一般的です。因子分析の結果として性格特性が得られる，などといったりするわけですから，これはもう心理学の王道中の王道，と言えるかもしれません。しかし実際は統計ソフトウェアが計算してくれて，何だかわからないまま使っている，という人も少なくありません。質問紙調査で分析して何かがわかる，というのはどういうことなのか，原理的なところから解説を始めていきます。またこの講義ではその基礎的なところ，数式レベルでしっかりと把握した上で，数値例に進みます。講義が終わる頃には，意味内容をしっかり理解したうえで因子分析を使えるようになっていることが期待されます。

この授業の後半(後期)は，実践・演習を主にした学習になります。我々が手にしたデータは，何らかのモデル・仮定のもとで生成されたものである，という考え方に立ち，\textbf{データから生成メカニズムを推測する}ことにチャレンジします。これは非常に夢のある話です。だってデータ生成メカニズムというのは，私たちの心の中の機序そのものだからです。推測に当たってはさまざまな前提・仮定が必要ですが，得られる結果は非常に含蓄に富み，さまざまな角度から心を考えさせてくれるヒントになります。線形モデルは直線的な関係でしたが，より柔軟で豊富な表現力を持つ技術へと理解を進めていきましょう。

今回は初回ですので，基礎的な心理統計で学んだことを改めて確認・復習することを中心に，今後の講義でも用いられる数学的記号の準備をします。

\section{正規線形モデルの世界}

因子分析法や(重)回帰分析では，基本的に扱う変数が複数あります。これまでの相関・回帰分析や，実験計画の中では，説明変数と被説明変数がひとつづつ，という二変数の世界でした。これが多くなったシーンは，一般に\keyterm{多変量解析}{たへんりょうかいせき}{Multivariate Analysis}と呼ばれます。変量あるいは\keyterm{変数}{へんすう}{Variables}は，ケースごとに変わる数のことを指します。変わる「数」と言っていますが，この後述べるように数字以外のものも数字として扱いますので，「ケースごとに変わるもの」の総称だと思ってください。多変量はそれが多くあるもの，データセット全体を指します。\textbf{たくさんのデータセットの中から意味ある情報を引きだすこと}，これが多変量解析の目的です。

イメージとしては，表計算ソフトのスプレッドシートに数字がずらっと並んでいる世界です。例えば性格検査の尺度の一種である，YG性格検査は被験者に120の項目について回答させます。ひとりにつき120の変数があり，これを何百，何千人に対して実施するのですから，非常に大量のデータになっているわけです。スプレッドシートにデータを入れていくときは一般に，一行(横方向)に1ケース(1オブザベーション，1個人)であるようにし，列方向(縦方向)に変数を並べるのが一般的です。数学記号では次のように表現します。

\[ x_{ij }\]

ここで$i$は個人のID

単変量ではなく多変量を扱う統計の領域に入るので，多変量データとはどのようなものであるかに言及した上で，本講義の扱う領域を概観する。心理統計の応用的分野では，正規分布を仮定した線形モデルがその大半を占めている。正規線形モデルに含まれるさまざまな下位モデルの名称を紹介するとともに，構造方程式モデリングに統合されることや，非線形なモデルとの違いについて理解する。加えて正規分布ではない分布を扱うモデルも増えてきている昨今，これらについてのモデリングアプローチの存在についても講義する。
    % \begin{flushright}
    %     正規線形モデルの枠組みについては，$\to$\citet{三中2018}参照。
    % \end{flushright}
    
\section{尺度の四水準}心理統計の基礎で触れたが，データ化として扱う数値はその尺度水準によってどのような計算が可能かということに違いが生じる。このことは，そのまま分析モデルや名称の違いに繋がるため，改めて名義，順序，間隔，比率の4水準を確認しておく。
    % \begin{flushright}
    %     $\to$ \citet{stevens1946}の論文は短く，ネットで読むこともできる。
    % \end{flushright}
    
\section{平均と分散}間隔尺度水準以上の数字であれば，平均値や分散，標準偏差によってその特徴を要約することができる。ここではこれらの代表地の表記について，数学的記号とともに確認する。加えて，分散式を展開して表現したものや，分散がデータから得られる情報の上限であることを確認する。

\section{共分散と相関係数}共分散やそれを標準化した相関係数は，複数の変数間関係を表現する最もシンプルなものの一つである。ここではこれらの複数の変数間に関わる代表値について，数学的記号とともに確認する。加えて，この他の関係の表現方法として，距離や共頻度などの共変量について解説し，それらの違いに応じて統計モデルが変わりうることを確認する。



\section{課題}
\paragraph{心理学とは何だろうか}
みなさんの中で心理学とはどういうイメージを持っていますか。この講義を受けるまえはどのように考えていましたか。この講義を受けた後で何か変わりましたか。できれば今考えている，あなたなりの心理学像を，どこかに書き起こしておくと良いでしょう。前期の終わり，あるいは一年の終わり，そして四年後にそれがどう変わるかを楽しみに。
\paragraph{心理学にどうして統計が必要なのか}
第1回は心理統計の中身に入る前に，どうして心理学なのに数字を扱って統計を学ぶのかについてを解説しました。みなさんの身近な人に，みなさんなりの言葉でこのことを説明してみてください。そのときに，心理統計と心理測定の違い，科学としての心理学，というキーワードを入れながら説明できれば完璧です\footnote{また後ほど触れることになりますが，ある人が何かを理解している，わかっているということは，そのことについて様々な表現を延々と生み出すことができる，ということでもあります。もし言葉に詰まったり，うまく説明できないことがあれば，そこの理解ができていないと考えられます。このことは第\ref{chapter5}講で触れる古典的テスト理論とも関係しますし，心理学の基本モデルである図\ref{fig::01_01}とも直結する観点です。}。
\end{document} 
